% =============================================================================
% ANEXO V — Mapa de Pantallas de la SPA
% =============================================================================
\chapter{Mapa de Pantallas de la SPA}
\label{cap:anexo_pantallas}

Inventario de las páginas de la SPA (\comando{EthosFront}), organizadas por rol y protegidas por
los \textit{guards} de enrutamiento (Anexo~\ref{cap:anexo_frontend_inv}).\par

\section{Pantallas por Rol}
\label{sec:anexo_pantallas_rol}
\renewcommand{\arraystretch}{1.3}
\fontsize{8.8}{10.5}\selectfont\sffamily
\noindent
\begin{tabularx}{\textwidth}{|l|X|}
\hline
\rowcolor{colorPrimario}
\textcolor{white}{\textbf{Grupo (\texttt{pages/})}} & \textcolor{white}{\textbf{Pantallas}} \\ \hline
\comando{auth/} & Login, registro, recuperación de contraseña, \textit{callback} OAuth. \\ \hline
\rowcolor{grisTecnico}
\comando{dashboard/} & Proyectos, experiencia, formación, skills, CV Studio, conexiones, portafolio, chat, preferencias. \\ \hline
\comando{recruiter/} & Talent Discovery (Kanban), dashboard analítico, leaderboard, chat, \textit{insights}. \\ \hline
\rowcolor{grisTecnico}
\comando{admin/} & Métricas, moderación, gestión de perfiles y skills, correos. \\ \hline
\comando{public/} & Landing, portafolio público (con/sin contraseña), explorar talento, páginas legales. \\ \hline
\end{tabularx}\normalfont\normalsize
\par\vspace{0.3cm}

\section{Protección de Rutas}
\label{sec:anexo_pantallas_guards}
\renewcommand{\arraystretch}{1.3}
\fontsize{9}{10.8}\selectfont\sffamily
\noindent
\begin{tabularx}{\textwidth}{|l|X|}
\hline
\rowcolor{colorPrimario}
\textcolor{white}{\textbf{Grupo}} & \textcolor{white}{\textbf{Protección}} \\ \hline
\comando{auth/}, \comando{public/} & Acceso libre (sin sesión). \\ \hline
\rowcolor{grisTecnico}
\comando{dashboard/} & \comando{ProtectedRoute} + rol PROFESSIONAL. \\ \hline
\comando{recruiter/} & \comando{ProtectedRoute} + rol RECRUITER. \\ \hline
\rowcolor{grisTecnico}
\comando{admin/} & \comando{ProtectedRoute} + rol ADMIN (dominio). \\ \hline
\end{tabularx}\normalfont\normalsize

\begin{NotaTecnica}
El portafolio público (\comando{public/}) admite protección con contraseña por configuración del
profesional; el \comando{RouteErrorBoundary} captura errores de renderizado sin romper la SPA.
\end{NotaTecnica}

\clearpage

% =============================================================================
\chapter{Recorridos de Usuario (User Journeys)}
\label{cap:anexo_journeys}

Recorridos clave de los actores principales, encadenando pantallas y endpoints reales.\par

\section{Profesional: del Registro al Portafolio Público}
\label{sec:anexo_journey_prof}
\begin{enumerate}[noitemsep]
    \item Registro en \comando{/api/auth/register} (rol PROFESSIONAL) y redirección al dashboard.
    \item Completar el perfil (\comando{/api/v1/profile}) y la trayectoria (experiencia, formación).
    \item Añadir hard/soft skills y crear proyectos con evidencias (subida Zero Trust).
    \item Editar el CV en CV Studio y compilarlo a PDF.
    \item Configurar visibilidad y SEO; publicar el portafolio (opcionalmente con contraseña).
\end{enumerate}

\section{Reclutador: del Descubrimiento a la Contratación}
\label{sec:anexo_journey_recruiter}
\begin{enumerate}[noitemsep]
    \item Registro (rol RECRUITER) y acceso a Talent Discovery.
    \item Búsqueda NLP y exploración de talento; \textit{like} a candidatos.
    \item Organizar candidatos en el pipeline Kanban (\comando{saved}$\to$\comando{hired}).
    \item Tomar notas privadas y consultar \textit{insights} de IA y dashboards.
    \item Contactar por chat en tiempo real y avanzar la etapa hasta \comando{hired}.
\end{enumerate}

\section{Administrador: Moderación y Métricas}
\label{sec:anexo_journey_admin}
\begin{enumerate}[noitemsep]
    \item Acceso al panel (\comando{/api/admin/metrics/overview}).
    \item Revisión de métricas de crecimiento, roles y \textit{top skills}.
    \item Moderación de perfiles (soft-delete y restauración).
    \item Gestión del catálogo de skills y envío de correos.
\end{enumerate}

\clearpage
